\documentclass[10pt,a4paper]{amsart}
\usepackage[margin=19mm]{geometry}
\usepackage{amsmath,amssymb,mathtools}
\usepackage[scr=rsfs,cal=euler]{mathalfa}
\usepackage{xfrac}
\usepackage[unicode,hidelinks,bookmarks=false]{hyperref}
\newcommand{\ie}{i.e.\ }
\newcommand{\cf}{cf.\ }
\newcommand{\resp}{respectively\ }
\newcommand{\Subsection}{\subsection}
\providecommand{\nobreakdash}{\nobreak\hbox{-}}
\setlength{\emergencystretch}{2em}
\multlinegap0pt
\usepackage{bm}
\usepackage{bbm}
\usepackage{dsfont}
\usepackage{xspace}
\usepackage{cancel}
\usepackage{mathtools}
\usepackage[shortlabels]{enumitem}
\usepackage{amsthm}
\makeatletter
\@ifundefined{theorem}{\newtheorem{theorem}{Theorem}}{}
\@ifundefined{lemma}{\newtheorem{lemma}[theorem]{Lemma}}{}
\makeatother
\def\sol{\mathrm{sol}} 
\title[The solvable Graph of a finite-dimensional Lie Algebra]{The solvable Graph of a finite-dimensional Lie Algebra}
\author{David Towers, Ismael Gutierrez, Luis Fernandez}
\date{Source: arXiv:2511.08290v1 (CC BY 4.0)}
\begin{document}
\maketitle

\noindent\emph{Source and licence.} The solvable Graph of a finite-dimensional Lie Algebra. Version arXiv:2511.08290v1 (CC BY 4.0), distributed under the Creative Commons Attribution 4.0 International licence, as recorded in the arXiv licence metadata for this exact version and shown on the abstract page https://arxiv.org/abs/2511.08290v1. Full rights notice: licenses/arxiv-2511.08290-CC-BY-4.0.txt.\par\smallskip

\noindent\textit{Editorial scope.} Counted unit: the Lemma whose source label is \texttt{solvabilizer1} in the article (source0.tex lines 419--426, source label \texttt{solvabilizer1}), delivered together with the article's own complete original proof of it (source0.tex lines 428--453, one \texttt{proof} environment of 26 lines). No mathematics is written, repaired or re-derived: statement and proof are the article's own text line for line. Required context retained uncounted: none. Every definition, display and label of those lines is retained unchanged. Omitted from this package and not redistributed: 1--418, 427--427, 454--1353. Nothing inside a retained excerpt is omitted or altered: every retained body is delivered exactly as the article's lines are written, so no omission receipt of any kind accompanies this package. Citations: the retained excerpts carry no citation command at all, so no bibliography entry of the article is transcribed and no reference is redistributed. Reference adaptation: none was needed and none was made; every reference inside the delivered unit resolves inside it, so no unresolved reference or citation is produced. Printed slips kept literal: no printed word, symbol, hypothesis or number of the counted statement or of its proof was altered. Layout and class: the article's own class is \texttt{amsart} with packages including amsmath, amsthm, amscd, amssymb, amsfonts, color, graphicx, mathrsfs, figsize, algorithm, algpseudocode, enumerate, float, tikz; the wrapper is the lane's pinned \texttt{amsart} editorial wrapper at 10pt on A4, which restates the theorem-like environments the retained text needs and adds the article's own macro definitions that the retained text needs (\texttt{\textbackslash sol}), with extra wrapper packages (bm, bbm, dsfont, xspace, cancel, mathtools, enumitem). The delivered numbering is the unit's own. Assets: no retained span contains a figure environment, a graphics-inclusion command, a PDF-page inclusion, a table or a TikZ picture; no image file is copied into this package, the package declares no asset dependency and no image file is redistributed with it. Licence: version arXiv:2511.08290v1 of this article is distributed under the Creative Commons Attribution 4.0 International licence, as recorded in the arXiv licence metadata for this exact version and shown on the abstract page https://arxiv.org/abs/2511.08290v1. A full rights notice is delivered at licenses/arxiv-2511.08290-CC-BY-4.0.txt. Characters: every retained excerpt is pure ASCII, so the delivered engine, XeLaTeX with its default Latin Modern text font, reports no missing character.
\begin{lemma}\label{solvabilizer1}
Let $L$ be a Lie algebra of dimension $n$ over a field $F$ with $|F|> n-2$, and let $x\in L$. Then the following statements are equivalent:
\begin{enumerate}[\rm(i)]
  \item $\sol_L(x)$ is a \emph{subspace} of $L$;
  \item $\sol_L(x)$ is a \emph{maximal solvable subalgebra} of $L$;
  \item $x$ lies in a unique maximal solvable subalgebra of $L$.
\end{enumerate}
\end{lemma}

\begin{proof}
First, we claim that
\begin{equation}\label{sol1}
\sol_L(x) = \bigcup\big\{M\le L \mid M \text{ maximal solvable and } x\in M \big\}.
\end{equation}
In fact, let $y\in \sol_L(x)$. Then $\langle x,y\rangle$ is solvable, hence it is contained in some
maximal solvable subalgebra $M$ with $x,y\in M$; thus $y$ lies in the right side of \eqref{sol1}.

Conversely, if $y$ belongs to a maximal solvable subalgebra $M$ containing $x$, the subalgebra $\langle x, y\rangle\subseteq M$ is solvable, so $y\in\sol_L(x)$. This proves \eqref{sol1}
\medskip

\noindent\rm{(ii) $\Rightarrow$ (i).} Any maximal solvable subalgebra is of course a subspace.
\medskip

\noindent\rm{(iii) $\Rightarrow$ (ii).} If there is a unique maximal solvable subalgebra $M$
containing $x$, then by \eqref{sol1} we have $\sol_L(x)=M$, which is maximal solvable by hypothesis.
\medskip

\noindent\rm{(ii) $\Rightarrow$ (iii).} Suppose $M:=\sol_L(x)$ is a maximal solvable subalgebra of $L$. Let now $N$ be another maximal solvable subalgebra $N\neq M$ contained $x$, then $N \subseteq \sol_L(x) = M$ by \eqref{sol1}; maximality forces
$N=M$, a contradiction.  Hence $x$ lies in exactly one maximal
solvable subalgebra.
\medskip


\noindent\rm{(i) $\Rightarrow$ (iii).}  Assume $\sol_L(x)$ is a subspace of $L$.  Suppose further that $x$ lies in $k>1$ maximal solvable subalgebras $M_1, \ldots, M_k$ with $M_i \not\subseteq \cup_{j\neq i} M_j$. Since $\dim M_i\geq 2$, $n\geq \dim \cup_{i=1}^kM_i\geq k+1$. By Lemma \ref{solvabilizer1}, $|F|\leq k-1\leq n-2$, a contradiction. Therefore, there is exactly one maximal solvable subalgebra containing $x$.
\end{proof}

\end{document}
